% Danh muc tai lieu, chep tu paper/paper1/bibliography.tex (bo danh dau ban sua TETC).
% Thu tu muc = thu tu trich dau tien trong bai (sap lai 20-11).
% Van dang kieu IEEE (thebibliography). Doi sang .bib + sn-mathphys-num neu QMI doi:
% xem docs/PAPER1_KE_HOACH_QMI.md, Giai doan 2.
\begin{thebibliography}{36}

\bibitem{ref24}
A. Khraisat, I. Gondal, P. Vamplew, and J. Kamruzzaman, ``Survey of intrusion detection systems: techniques, datasets and challenges,'' Cybersecurity, vol. 2, no. 1, p. 20, 2019.

\bibitem{ref1}
C. Yin, Y. Zhu, J. Fei, and X. He, ``A deep learning approach for intrusion detection using recurrent neural networks,'' IEEE Access, vol. 5, pp. 21954--21961, 2017.

\bibitem{ref2}
R. Vinayakumar, M. Alazab, K. P. Soman, P. Poornachandran, A. Al-Nemrat, and S. Venkatraman, ``Deep learning approach for intelligent intrusion detection system,'' IEEE Access, vol. 7, pp. 41525--41550, 2019.

\bibitem{ref3}
H. Liu and B. Lang, ``Machine learning and deep learning methods for intrusion detection systems: a survey,'' Applied Sciences, vol. 9, no. 20, p. 4396, 2019.

\bibitem{havlicek2019}
V. Havl\'i\v{c}ek, A. D. C\'orcoles, K. Temme, A. W. Harrow, A. Kandala, J. M. Chow, and J. M. Gambetta, ``Supervised learning with quantum-enhanced feature spaces,'' Nature, vol. 567, no. 7747, pp. 209--212, 2019.

\bibitem{schuld2019}
M. Schuld and N. Killoran, ``Quantum machine learning in feature Hilbert spaces,'' Physical Review Letters, vol. 122, no. 4, p. 040504, 2019.

\bibitem{ref4}
M. Tavallaee, E. Bagheri, W. Lu, and A. A. Ghorbani, ``A detailed analysis of the KDD CUP 99 data set,'' in Proc. IEEE Symp. Computational Intelligence for Security and Defense Applications, 2009, pp. 1--6.

\bibitem{ref25}
N. Moustafa and J. Slay, ``UNSW-NB15: a comprehensive data set for network intrusion detection systems,'' in Military Communications and Information Systems Conference (MilCIS), 2015, pp. 1--6.

% [26] cua ban TETC (Rahman et al., IEEE Access) da go: khong xac minh duoc
% su ton tai; thay bang kaissar2026.

\bibitem{ref13}
M. O. Kalinin and V. M. Krundyshev, ``Security intrusion detection using quantum machine learning techniques,'' Journal of Computer Virology and Hacking Techniques, vol. 19, no. 1, pp. 125--136, 2023.

\bibitem{ref14}
H. Suryotrisongko and Y. Musashi, ``Evaluating hybrid quantum-classical deep learning for cybersecurity botnet DGA detection,'' Procedia Computer Science, vol. 197, pp. 223--229, 2022.

% SUA theo R2: so hieu bai 116990F -> 116990B (xac minh qua SPIE Digital Library va NASA ADS).

\bibitem{ref15}
E. D. Payares and J. C. Mart\'inez-Santos, ``Quantum machine learning for intrusion detection of distributed denial of service attacks: a comparative overview,'' in Proceedings of SPIE: Quantum Computing, Communication, and Simulation, vol. 11699, 2021, p. 116990B.

\bibitem{kaissar2026}
A.~Kaissar, A.~B. Nassif, and A.~Bouridane, ``Enhancing network intrusion
detection with quantum machine learning: A comprehensive survey of methods,
metrics, and applications,'' \emph{Future Internet}, vol.~18, no.~5,
art. no. 234, 2026, doi: 10.3390/fi18050234.
% Thay cho [26] cua ban da nop. Xac minh doc lap: DOI phan giai dung bai,
% MDPI ghi Future Internet 18(5):234, 2026.

\bibitem{thanasilp2024}
S. Thanasilp, S. Wang, M. Cerezo, and Z. Holmes, ``Exponential concentration in quantum kernel methods,'' Nature Communications, vol. 15, no. 1, p. 5200, 2024.

\bibitem{ref5}
J. Biamonte, P. Wittek, N. Pancotti, P. Rebentrost, N. Wiebe, and S. Lloyd, ``Quantum machine learning,'' Nature, vol. 549, no. 7671, pp. 195--202, 2017.

\bibitem{ref6}
M. Cerezo, G. Verdon, H.-Y. Huang, L. Cincio, and P. J. Coles, ``Challenges and opportunities in quantum machine learning,'' Nature Computational Science, vol. 2, no. 9, pp. 567--576, 2022.

\bibitem{ref7}
P. Rebentrost, M. Mohseni, and S. Lloyd, ``Quantum support vector machine for big data classification,'' Physical Review Letters, vol. 113, no. 13, p. 130503, 2014.

\bibitem{ref9}
M. Schuld, ``Supervised quantum machine learning models are kernel methods,'' arXiv preprint arXiv:2101.11020, 2021.

\bibitem{ref18}
Y. Liu, S. Arunachalam, and K. Temme, ``A rigorous and robust quantum speed-up in supervised machine learning,'' Nature Physics, vol. 17, no. 9, pp. 1013--1017, 2021.

\bibitem{ref19}
H.-Y. Huang, M. Broughton, M. Mohseni, R. Babbush, S. Boixo, H. Neven, and J. R. McClean, ``Power of data in quantum machine learning,'' Nature Communications, vol. 12, no. 1, p. 2631, 2021.

\bibitem{ref20}
J. M. K\"ubler, S. Buchholz, and B. Sch\"olkopf, ``The inductive bias of quantum kernels,'' Advances in Neural Information Processing Systems, vol. 34, pp. 12661--12673, 2021.

\bibitem{cristianini2002}
N. Cristianini, J. Shawe-Taylor, A. Elisseeff, and J. Kandola, ``On kernel-target alignment,'' Advances in Neural Information Processing Systems, vol. 14, 2002.

\bibitem{cortes2012}
C. Cortes, M. Mohri, and A. Rostamizadeh, ``Algorithms for learning kernels based on centered alignment,'' Journal of Machine Learning Research, vol. 13, pp. 795--828, 2012.

\bibitem{ref10}
J. R. McClean, S. Boixo, V. N. Smelyanskiy, R. Babbush, and H. Neven, ``Barren plateaus in quantum neural network training landscapes,'' Nature Communications, vol. 9, no. 1, p. 4812, 2018.

\bibitem{ref21}
S. Sim, P. D. Johnson, and A. Aspuru-Guzik, ``Expressibility and entangling capability of parameterized quantum circuits for hybrid quantum-classical algorithms,'' Advanced Quantum Technologies, vol. 2, no. 12, p. 1900070, 2019.

\bibitem{qmi2026}
F.~Cirillo, C.~Esposito, and J.~T. Seo, ``Benchmarking quantum machine
learning methods for intrusion detection on noisy quantum computers,''
\emph{Quantum Mach. Intell.}, vol.~8, no.~1, art. no.~27, 2026,
doi: 10.1007/s42484-026-00379-4.
% Tac giả lay tu CrossRef API theo DOI (api.crossref.org/works/<doi>) roi
% kiem cheo doc lap bang tim kiem -- khop ca ba ten, volume, so bai, nam.
% Trang Springer chan toan van, nhung metadata thi khong chan: lan truoc
% bo cuoc qua som o trang HTML.

\bibitem{bowles2024}
J.~Bowles, S.~Ahmed, and M.~Schuld, ``Better than classical? The subtle art of
benchmarking quantum machine learning models,'' 2024, \emph{arXiv:2403.07059}.

\bibitem{schnabel2024}
J.~Schnabel and M.~Roth, ``Quantum kernel methods under scrutiny: A benchmarking
study,'' 2024, \emph{arXiv:2409.04406}.

\bibitem{gillani2026}
S.~A. Gillani, M.~S.~A. Baig, S.~M. Shah, A.~Ali, and H.~Siddiqui, ``How quantum
is the advantage? A fair, calibration- and noise-aware benchmark and attribution
audit of quantum machine learning for network intrusion detection,'' 2026,
\emph{arXiv:2608.18155}.

\bibitem{cortes1995}
C. Cortes and V. Vapnik, ``Support-vector networks,'' Machine Learning, vol. 20, no. 3, pp. 273--297, 1995.

\bibitem{scholkopf2002}
B. Sch\"olkopf and A. J. Smola, Learning with Kernels: Support Vector Machines, Regularization, Optimization, and Beyond. MIT Press, 2002.

\bibitem{ref32}
J. Cohen, Statistical Power Analysis for the Behavioral Sciences, 2nd ed. Hillsdale, NJ, USA: Lawrence Erlbaum Associates, 1988.

\bibitem{ref31}
T. G. Dietterich, ``Approximate statistical tests for comparing supervised classification learning algorithms,'' Neural Computation, vol. 10, no. 7, pp. 1895--1923, 1998.

\bibitem{ref30}
F. Pedregosa et al., ``Scikit-learn: machine learning in Python,'' Journal of Machine Learning Research, vol. 12, pp. 2825--2830, 2011.

\bibitem{ref29}
A.~Javadi-Abhari et al., ``Quantum computing with Qiskit,'' 2024, \emph{arXiv:2405.08810}, doi: 10.48550/arXiv.2405.08810.
% Doi 20-11 (chu nhan gat): ban cu chi ghi "Qiskit Contributors, 2024". Muc nay
% lay tu CITATION.bib cua repo Qiskit/qiskit (cach trich Qiskit tu de nghi).

\bibitem{ref33}
J. Qui\~nonero-Candela, M. Sugiyama, A. Schwaighofer, and N. D. Lawrence, Eds., Dataset Shift in Machine Learning. Cambridge, MA, USA: MIT Press, 2009.

% Da go 20-11 (khong muc nao trich): carducci2026 (chi doc duoc tom tat, bo
% muc Carducci o Related); ref8 Preskill NISQ (bai moi khong dung chu NISQ);
% ref27 trang web IBM roadmap (muc nhieu da trich Qiskit, ref29).
\end{thebibliography}
